\documentclass[11pt]{amsart}

\usepackage{amsmath,amssymb,amsthm}
\usepackage{stmaryrd}     % \llbracket \rrbracket
\usepackage{mathtools}
\usepackage{booktabs}
\usepackage{array}
\usepackage[hidelinks]{hyperref}
\usepackage{enumitem}

% ---- theorem environments ----
\theoremstyle{plain}
\newtheorem{theorem}{Theorem}[section]
\newtheorem{lemma}[theorem]{Lemma}
\newtheorem{proposition}[theorem]{Proposition}
\newtheorem{corollary}[theorem]{Corollary}
\newtheorem{conjecture}[theorem]{Conjecture}
\theoremstyle{definition}
\newtheorem{definition}[theorem]{Definition}
\newtheorem{remark}[theorem]{Remark}
\newtheorem{example}[theorem]{Example}

% ---- macros ----
\newcommand{\C}{\mathcal{C}}
\newcommand{\D}{\mathcal{D}}
\newcommand{\Fam}{\mathrm{Fam}}
\newcommand{\ext}[1]{\llbracket #1 \rrbracket}
\newcommand{\Setb}{\mathbf{Set}}
\newcommand{\Vecb}{\mathbf{Vec}}
\newcommand{\Vecfd}{\mathbf{Vec}_{\mathrm{fd}}}
\newcommand{\Gl}{\mathrm{Gl}}
\newcommand{\op}{\mathrm{op}}
\newcommand{\Hom}{\mathrm{Hom}}
\newcommand{\comp}{\mathbin{\triangleleft}}
\newcommand{\compDJN}{\mathbin{\triangleleft_{\mathrm{DJN}}}}
\newcommand{\one}{\mathbf{1}}
\newcommand{\gena}[1]{\langle #1 \rangle}
\newcommand{\pizero}{\pi_0}
\DeclareMathOperator{\supp}{supp}

\begin{document}

\title[Admissibility of the composition product over a base]
{When does substitution exist? \\ Admissibility of the composition product on $\Fam(\C^{\op})$}

\author{MacBeth}
\address{Kodamai}
\email{neil@kodamai.com}

\date{September 2026}

\begin{abstract}
A polynomial functor is a coproduct of representables, and substituting one into the variable
of another gives the composition product $\comp$ that organises polynomial functors into a
monoidal category. Fixing a closed symmetric monoidal base $\C$ and taking positions in $\C$
yields the external-shape category $\Fam(\C^{\op})$, with extension $\ext{S,P}X=\coprod_{s}[P_s,X]$;
one then asks for which bases $\C$ this category still carries $\comp$. We show the question is not
graded but nearly dichotomous. A single necessary condition -- the \emph{Shape Lemma}, that
$[P,\,T\cdot\one]$ must be a copower of the terminal object -- forces the answer on both extremes.
On any nontrivial infinitary-lextensive cartesian-closed base the mere existence of $\comp$ forces
the monoidal unit to be \emph{connected}, whence the extension is fully faithful, substitution is
canonical, and the slot-reindexing left adjoint $F_q\dashv({-})\comp q$ exists for every $q$. On this
pole we obtain a complete characterisation: $\Fam(\C^{\op})$ carries $\comp$ if and only if the
connected-components functor $\pizero$ preserves finite products. This one criterion unifies the two
ways admissibility can fail -- a disconnected unit (as for $\Setb\times\Setb$) and, more subtly, a
connected unit whose $\pizero$ is still not multiplicative. The latter is realised by the Artin
gluing $\Gl((-)^2)$, a Grothendieck topos with a \emph{connected} unit that is nonetheless
inadmissible because $\pizero(K\times K)=2\neq1$ for the single edge $K$ (Weichsel's theorem); so
connectedness of the unit is necessary but not sufficient. On an additive base substitution collapses
to the Dirichlet tensor and the unit is always disconnected. Between the poles the admissible region
is thin but non-empty: $\Setb\times\Vecfd$ is admissible with a disconnected unit, so the
extensive-pole characterisation does not extend to a global one. The organising fact is that when the
unit is connected the extension is injective on objects, so $\comp$ is determined by $\C$; when it is
disconnected $\comp$ is a choice, and the base-change question stops being about $\C$ at all.
\end{abstract}

\maketitle

\section{Introduction}

\subsection*{The problem}
A polynomial functor $\Setb\to\Setb$ is a coproduct of representables,
\[
   \ext{S,P}\,X \;=\; \coprod_{s\in S} X^{P_s},
\]
indexed by a set $S$ of \emph{shapes} with a set $P_s$ of \emph{positions} at each shape.
Two polynomial functors compose to a third, and the resulting operation -- plugging one polynomial
into the variable of another -- is the \emph{composition product} $\comp$, characterised by
$\ext{p\comp q}\cong\ext p\circ\ext q$. It makes polynomial functors the objects of a monoidal
category whose comonoids are exactly small categories \cite{NiuSpivak,DJN}, and it is the algebraic
heart of the ``polynomial'' account of interaction, dynamics, and data.

Nothing in the definition of a polynomial functor needs the exponents to be indexed by \emph{sets}.
Fix instead a closed symmetric monoidal cocomplete base $(\C,\otimes,I,[-,-])$. An object of
$\Fam(\C^{\op})$ is a set $S$ of shapes together with \emph{position objects} $P_s\in\C$, and its
extension uses the internal hom:
\[
   \ext{S,P}\,X \;=\; \coprod_{s\in S}[P_s,X]\;:\;\C\to\C .
\]
For $\C=\Setb$ this is the classical construction and $\ext{-}$ is fully faithful -- the
Abbott--Altenkirch--Ghani representation theorem \cite{AAG}. For $\C=\Vecb$ it is the category of
linear containers. Dorta, Jarvis and Niu \cite{DJN} study a closely related \emph{generalised
polynomial} category over a base and build a Dirichlet-style tensor and a (weighted) composition
product on it; they explicitly leave conditions on the base for future work
(\cite[\S6]{DJN}). This paper answers, for the external-shape category $\Fam(\C^{\op})$, the most
basic such question:
\begin{center}
\emph{for which bases $\C$ does $\Fam(\C^{\op})$ carry the composition product $\comp$ at all?}
\end{center}

\subsection*{The gap: the question has the wrong shape}
It is tempting to expect a graded answer, as one gets for the Dirichlet tensor
$p\otimes q=(S\times T,\,P_s\otimes Q_t)$, which exists over \emph{every} closed monoidal base and
whose finer properties (closedness, and so on) degrade continuously as the base varies. The
composition product is not like this. It \emph{mostly does not exist}, and where it does exist its
existence is not a matter of degree: on the natural ``set-like'' bases the mere availability of
$\comp$ already forces a rigid package of further structure, while on ``linear'' bases $\comp$
degenerates into $\otimes$ and the interesting adjoints vanish. The right output is therefore not a
gradation but a \emph{map} of the space of bases into a small number of qualitatively distinct
regions -- and the boundaries between them are sharp theorems, not thresholds.

\subsection*{The results}
We call $\C$ \emph{$\comp$-admissible} if the image of $\ext{-}$ is closed under composition:
for all $p,q$ there is an object $p\comp q$ with $\ext{p\comp q}\cong\ext p\circ\ext q$
(Definition~\ref{def:admissible}). This is the weakest possible reading of ``$\comp$ exists'', so
every negative result below is as strong as it can be. Everything is controlled by one necessary
condition.

\begin{theorem}[Shape Lemma, Lemma~\ref{lem:shape}]
If $\C$ is $\comp$-admissible then for every object $P$ and every set $T$ the internal hom
$[P,\;T\cdot\one]$ is a copower of the terminal object $\one$, where $T\cdot\one=\coprod_T\one$.
\end{theorem}

The terminal object $\one=[0,X]$ always exists (Lemma~\ref{lem:terminal}), and the Shape Lemma says
the base must be able to \emph{see} the external shape set inside itself, as a copower of $\one$.
Over $\Setb$ this is vacuous ($\one=1$ generates); over $\Vecb$ it is vacuous the other way
($\one=0$, so every copower is $0$ and the shape data is invisible). Between the degeneracies it has
teeth, and it dictates the whole map.

\emph{The extensive pole is rigid.} On the set-like extreme, admissibility is not just consistent
with the rest of the structure -- it \emph{forces} it.

\begin{theorem}[Extensive rigidity, Theorem~\ref{thm:B} and Theorem~\ref{thm:leftadj}]
\label{intro:B}
Let $\C$ be nontrivial, infinitary lextensive and cartesian closed. If $\C$ is $\comp$-admissible
then its unit $\one$ is \emph{connected} (i.e.\ $\C(\one,-)$ preserves small coproducts).
Consequently the extension $\ext{-}$ is fully faithful, and the reindexing functor
$F_q=\Fam(\ext q^{\op})$ is left adjoint to $({-})\comp q$ for \emph{every} $q$.
\end{theorem}

So on this pole admissibility hands over the whole package at once: a connected unit, a fully faithful
extension, and the reindexing left adjoint for every $q$. In particular $\Setb\times\Setb$ is
inadmissible, because its unit $(1,1)$ is disconnected. The converse fails, however -- connectedness
alone does \emph{not} suffice -- and pinning down exactly which extensive bases are admissible is the
next result.

\emph{The extensive pole is completely characterised.} Rigidity extracts connectedness from
admissibility but does not run backwards. The exact condition is that the Shape Lemma, necessary in
general, becomes sufficient -- equivalently, that the connected-components functor is multiplicative.

\begin{theorem}[Characterisation, Corollary~\ref{cor:char}]\label{intro:char}
Let $\C$ be nontrivial, infinitary lextensive and cartesian closed. Then $\C$ is $\comp$-admissible if
and only if the Shape Lemma holds, i.e.\ $[P,\,T\cdot\one]$ is a copower of $\one$ for all $P$ and all
$T$. When $\C$ admits a connected-components functor $\pizero\dashv(-)\cdot\one$ (as every Grothendieck
topos does) this reads: $\C$ is $\comp$-admissible if and only if $\pizero$ preserves finite products.
\end{theorem}

The single criterion $\pizero(X\times Y)\cong\pizero X\times\pizero Y$ subsumes both obstructions
below: the disconnected unit of $\Setb\times\Setb$ and the connected-unit-but-non-multiplicative
gluing $\Gl((-)^2)$ are one phenomenon, not two.

\emph{The additive pole is flexible.} On the linear extreme, every object is
\emph{copower-tiny} ($[P,-]$ preserves coproducts), $\comp$ collapses to $\otimes$, and:

\begin{theorem}[Collapse, Lemma~\ref{lem:D} and Theorem~\ref{thm:2}]\label{intro:D}
If every object of $\C$ is copower-tiny then the unit is disconnected, and $({-})\comp q$ has a left
adjoint if and only if $q$ has a single shape.
\end{theorem}

\emph{The middle is thin but inhabited.} The characterisation above is confined to the extensive
pole; it does not extend to a global one, because there is an admissible base off the pole, with a
disconnected unit, that is neither extreme.

\begin{theorem}[Theorem~\ref{thm:setxvec}]\label{intro:setxvec}
The base $\Setb\times\Vecfd$ is $\comp$-admissible on its natural locus, non-collapse, non-cartesian,
and has a disconnected unit. Hence the hypotheses of Extensive rigidity are essential:
admissibility alone does not force connectedness.
\end{theorem}

Its admissibility comes from combining a \emph{rigid} distributive factor ($\Setb$) with a
\emph{flexible} copower-tiny factor ($\Vecfd$); admissibility is a per-object property
(Proposition~\ref{prop:absorptive}), and the two factors supply the two mechanisms.

\emph{Connected is not enough.} The sharpest boundary phenomenon is that connectedness of the unit,
which Extensive rigidity extracts \emph{from} admissibility, does not conversely \emph{imply} it.

\begin{theorem}[Gluing obstruction, Theorem~\ref{thm:G}]\label{intro:G}
The Artin gluing $\Gl((-)^2)$ is a Grothendieck topos with a connected unit that is nonetheless
$\comp$-inadmissible. The obstruction is that its connected-components functor $\pizero$ is not
multiplicative: for the single edge $K$ one has $\pizero(K\times K)=2\neq 1$ (Weichsel's theorem).
\end{theorem}

\emph{The organising fact.} Behind the whole map is one dichotomy about the extension itself.

\begin{theorem}[Determinacy, Theorem~\ref{thm:determinacy}]\label{intro:determinacy}
If $\one$ is connected then $\ext{-}$ is injective on objects up to isomorphism, so $p\comp q$ is
\emph{determined} by $\C$ and ``$({-})\comp q$ has a left adjoint'' is a property of $\C$. If $\one$
is disconnected this can fail: over $\Vecfd$ the objects $(\{*\},k^2)$ and $(\{1,2\},k)$ present the
same functor $X\mapsto X\oplus X$ but are not isomorphic, so $\comp$ is a \emph{choice}.
\end{theorem}

Determinacy is why the base-change question changes character across the map: exactly when the unit
is disconnected, $\comp$ stops being pinned down by $\C$, so ``the'' composition product -- and hence
its adjoints -- become properties of a choice rather than of the base.

The whole map is summarised in Table~\ref{tab:map}.

\begin{table}[h]
\renewcommand{\arraystretch}{1.25}
\setlength{\tabcolsep}{4pt}
{\footnotesize
\begin{tabular}{@{}lcccl@{}}
\toprule
base $\C$ & adm.? & unit conn.? & $\pizero$ mult.? & verdict / mechanism \\
\midrule
$\Setb$, conn.\ topos & yes & yes & yes & always (Thm.~\ref{thm:leftadj})\\
$\Vecfd$, additive & on locus & no (forced) & --- & iff $|T|=1$ (Thm.~\ref{thm:2})\\
$\Setb\times\Vecfd$ & yes (locus) & no & --- & rigid $\times$ flexible (Thm.~\ref{thm:B})\\
$\Setb\times\Setb$ & no & no & no & $\pizero$ non-mult (Cor.~\ref{cor:unif})\\
$\Setb_*$ (pointed) & no & no & --- & off pole; degree $a=-4$ (Thm.~\ref{thm:A})\\
$\Gl((-)^2)$ (gluing) & no & \textbf{yes} & \textbf{no} & $\pizero$ non-mult (Thm.~\ref{thm:G})\\
\bottomrule
\end{tabular}
}
\medskip
\caption{The map of bases. On the extensive pole (rows with a $\pizero$ column) admissibility
\emph{is} $\pizero$-multiplicativity (Corollary~\ref{cor:char}); the last two such rows share one
mechanism despite differing on connectedness. A dash marks a base off the extensive pole, where
$\pizero$ need not exist. The gluing row is the crux: a connected unit does not suffice.}
\label{tab:map}
\end{table}

\subsection*{Method}
Two tools do all the work. The first is the Shape Lemma, obtained by evaluating any candidate
$p\comp q$ at the terminal object. The second is the \emph{probe method}: a single comparison map
\[
   \gamma_{(X_d)}\;:\;\coprod_{d}\C(\one,X_d)\;\longrightarrow\;\C\!\Big(\one,\coprod_d X_d\Big),
   \qquad (d,f)\mapsto \iota_d\circ f,
\]
controls fullness of $\ext{-}$, existence of the left adjoint, and (through variants that replace the
probe $\one$ by another object) the collapse verdict. Connectedness of the unit is exactly
``$\gamma$ is always a bijection'', and this one map is why fullness of the extension and existence of
the reindexing adjoint are, as we note in Remark~\ref{rem:onelemma}, the same lemma applied twice.

\subsection*{Scope}
Throughout, $\Fam(\C^{\op})$ has a \emph{single external} shape set, with positions in $\C$; this is
the setting of \cite{AAG} and of \cite{DJN}. The distinction matters: internal (Gambino--Kock)
polynomial functors in a topos are always closed under composition \cite{GambinoKock}, and the
Gluing obstruction (Theorem~\ref{intro:G}) is precisely the gap between the internal composition and
the external-shape one, localised to a base where it bites. We use $\comp$ throughout for the
composition-representing product $\ext{p\comp q}\cong\ext p\circ\ext q$; this is \emph{not}
the weighted product $\compDJN$ of \cite{DJN}, which carries an extra factor and agrees with $\comp$
only when $\C=1$ (Remark~\ref{rem:djn}). Their Dirichlet tensor, however, does coincide with ours.

\subsection*{Outline}
Section~\ref{sec:prelim} fixes the category $\Fam(\C^{\op})$ and the two invariants (connected unit,
copower-tiny object). Section~\ref{sec:shape} proves the Shape Lemma. Section~\ref{sec:extensive} is
the extensive pole (Extensive rigidity and the forced package). Section~\ref{sec:char} shows the
Shape Lemma is also sufficient there and derives the $\pizero$-multiplicativity characterisation.
Section~\ref{sec:collapse} is the
additive pole. Section~\ref{sec:inadmissible} collects the three inadmissible bases and the
``connected is not sufficient'' phenomenon. Section~\ref{sec:middle} exhibits the inhabited middle
region. Section~\ref{sec:determinacy} proves Determinacy and explains the shape of the whole map.
Section~\ref{sec:conclusion} states what remains open.

\subsection*{Provenance and status}
All numbered results are proved. The results about $\Setb\times\Vecfd$ (Section~\ref{sec:middle}) are
proved on an explicitly stated locus; the general additive case with infinite shape sets is open
(Remark~\ref{rem:gap1}). The topos-theoretic and graph-theoretic inputs to
Theorem~\ref{thm:G} -- that Artin gluing of a left-exact functor is a topos \cite{MacLaneMoerdijk,
CarboniJohnstone}, and Weichsel's theorem on the graph tensor product \cite{Weichsel} -- are
classical; the specific computations they feed ($\pizero(K\times K)=2$, the explicit internal hom
$[K,2\cdot\one]$) were checked independently. The external distributive law and the composite formula
underlying Section~\ref{sec:char} (Theorem~\ref{thm:Dsuff}, Theorem~\ref{thm:char}) were verified as
natural isomorphisms by exhaustive computation over finite sets, with no failures.
Conjecture~\ref{conj:decomp} is flagged as such.

\section{Preliminaries}\label{sec:prelim}

Fix a closed symmetric monoidal category $(\C,\otimes,I,[-,-])$ with small coproducts $\coprod$;
write $0$ for the initial object. Closedness makes each $[-,X]$ a right adjoint, hence
$[-,X]$ sends coproducts to products and $Y\otimes-$ preserves the initial object.

\begin{definition}[Containers over $\C$]
$\Fam(\C^{\op})$ has objects $p=(S,P)$ with $S$ a small set of \emph{shapes} and $P_s\in\C$ a
\emph{position object} for each $s\in S$; a morphism $(S,P)\to(T,Q)$ is a function $f:S\to T$
together with maps $\varphi_s:Q_{f(s)}\to P_s$ in $\C$. Equivalently
\begin{equation}\label{eq:hom}
   \Fam(\C^{\op})\big((S,P),(T,Q)\big)\;=\;\prod_{s\in S}\ \coprod_{t\in T}\ \C(Q_t,P_s).
\end{equation}
We write $\gena{Z}:=(\{*\},Z)$ for the one-shape container on $Z$, so every object is
$\coprod_{s}\gena{P_s}$. The \emph{extension} is $\ext{S,P}X=\coprod_{s\in S}[P_s,X]$, a functor
$\Fam(\C^{\op})\to[\C,\C]$. The \emph{Dirichlet tensor} is $p\otimes q=(S\times T,\,P_s\otimes Q_t)$,
and $T\cdot X:=\coprod_{t\in T}X$ denotes the copower.
\end{definition}

For $\C=\Setb$ this is the usual container calculus, $\ext{S,P}X=\coprod_s X^{P_s}$; for $\C=\Vecb$,
$\ext{S,P}W=\bigoplus_s\Vecb(P_s,W)$.

\begin{definition}[Admissibility]\label{def:admissible}
$\C$ is \emph{$\comp$-admissible} if for all $p,q\in\Fam(\C^{\op})$ there is an object
$p\comp q\in\Fam(\C^{\op})$ with $\ext{p\comp q}\cong\ext p\circ\ext q$.
\end{definition}

\begin{lemma}[Terminal object; {\cite[Lemma~1.0]{proofconv}}]\label{lem:terminal}
$\C$ has a terminal object $\one:=[0,X]$ (independent of $X$), and $[Q,\one]\cong\one$ for all $Q$.
\end{lemma}
\begin{proof}
$\C(Y,[0,X])\cong\C(Y\otimes 0,X)\cong\C(0,X)=1$ since $Y\otimes-$ preserves $0$; and
$\C(Y,[Q,\one])\cong\C(Y\otimes Q,\one)=1$.
\end{proof}

\begin{lemma}[Isomorphisms in a free coproduct completion]\label{lem:iso}
$(S,P)\cong(S',P')$ in $\Fam(\D)$ if and only if there is a bijection $\beta:S\to S'$ with
$P_s\cong P'_{\beta s}$ for all $s$.
\end{lemma}
\begin{proof}
Mutually inverse morphisms are bijections on shapes with mutually inverse position legs; the converse
is clear from \eqref{eq:hom}.
\end{proof}

The two invariants of the base that govern everything are the following.

\begin{definition}[Connected unit; the comparison map $\gamma$]\label{def:connected}
For a small family $(X_d)_{d\in D}$ let
\[
   \gamma_{(X_d)}\;:\;\coprod_{d\in D}\C(\one,X_d)\;\longrightarrow\;\C\!\Big(\one,\coprod_{d\in D}X_d\Big),
   \qquad (d,f)\mapsto\iota_d\circ f .
\]
The unit $\one=I$ is \emph{connected} if $\gamma_{(X_d)}$ is a bijection for every small family, i.e.\
$\C(\one,-):\C\to\Setb$ preserves small coproducts.
\end{definition}

For $\C=\Setb$, $\C(1,-)=\mathrm{Id}$ and the unit is connected. For $\C=\Vecb$, $\C(k,-)$ is the
forgetful functor and $\gamma:\coprod_d V_d\to\bigoplus_d V_d$ is neither injective (the zero vector
sits in every summand) nor surjective (cross-summand sums are missed): the unit is disconnected. Any
base with a zero object has a disconnected unit.

\begin{definition}[Copower-tiny object]\label{def:tiny}
$P\in\C$ is \emph{copower-tiny} if $[P,-]$ preserves small coproducts, i.e.\ the canonical
$\coprod_t[P,X]\to[P,\,T\cdot X]$ is an isomorphism for all $X$ and all sets $T$.
\end{definition}

Over $\Vecfd$ every object is copower-tiny; over $\Setb$ the tiny objects are the subterminals.
The extension records the exact link between connectedness and fullness.

\begin{theorem}[{Fullness criterion, \cite[Thm.~2.1]{proofT1}}]\label{thm:T1}
For $\C$ closed symmetric monoidal with small coproducts, the extension
$\ext{-}:\Fam(\C^{\op})\to[\C,\C]$ is fully faithful if and only if the unit $\one$ is connected. In
particular the action of $\ext{-}$ on hom-sets is $\prod_s$ of the comparison map
$\gamma_{([Q_t,P_s])_t}$.
\end{theorem}

This is the general form of the Abbott--Altenkirch--Ghani theorem \cite{AAG}, whose hidden hypothesis
is that the terminal object of $\Setb$ is connected; extensivity is neither necessary nor sufficient
(see $\Setb\times\Setb$, Corollary~\ref{cor:setxset} below). We use only that connectedness gives
full faithfulness, hence injectivity on objects up to isomorphism.

\section{The Shape Lemma}\label{sec:shape}

Admissibility has exactly one elementary consequence, extracted by evaluating any candidate
composite at the terminal object.

\begin{lemma}[Shape Lemma]\label{lem:shape}
If $\C$ is $\comp$-admissible then for every $P\in\C$ and every small set $T$, the object
$[P,\,T\cdot\one]$ is a copower of $\one$.
\end{lemma}
\begin{proof}
Fix $T$ and take any $q=(T,Q)$ with shape set $T$ (e.g.\ $Q_t:=0$). By Lemma~\ref{lem:terminal},
$\ext q(\one)=\coprod_{t}[Q_t,\one]\cong T\cdot\one$. Put $p:=\gena{P}$, so
$\ext p(\ext q(\one))\cong[P,\,T\cdot\one]$. By admissibility there is $p\comp q=(D,N)$ with
$\ext{p\comp q}\cong\ext p\circ\ext q$; evaluating at $\one$ and using Lemma~\ref{lem:terminal}
again,
\[
   \ext{p\comp q}(\one)=\coprod_{d\in D}[N_d,\one]\cong\coprod_{d\in D}\one=D\cdot\one .
\]
Hence $[P,\,T\cdot\one]\cong D\cdot\one$.
\end{proof}

\begin{remark}[Why it has teeth only in the middle]\label{rem:teeth}
The shape set $D$ of $p\comp q$ is an \emph{external} set -- one set, not one per ``component'' of
$\C$ -- and the Shape Lemma says the base must realise it internally as a copower of $\one$. Over
$\Setb$ this is vacuous: $\one=1$ generates and every set is a copower of it, the copower count being
the decoration set $T^{P}$. Over $\Vecb$ it is vacuous for the opposite reason: $\one=0$ so
$E\cdot\one=0$ for every $E$, and the shape data is invisible to $\ext{-}$ -- which is exactly the
failure of faithfulness in Theorem~\ref{thm:T1}. Between the two degeneracies the Shape Lemma is a
genuine, falsifiable demand, and Sections~\ref{sec:extensive} and \ref{sec:inadmissible} are its two
edges.
\end{remark}

\section{The extensive pole: admissibility forces the package}\label{sec:extensive}

Throughout this section $\C$ is \emph{nontrivial} ($\one\not\cong 0$), \emph{infinitary lextensive}
(finite limits, and $\C/\!\coprod_d A_d\simeq\prod_d\C/A_d$ for all small coproducts) and
\emph{cartesian closed}, so $\otimes=\times$, $I=\one$, $[X,Y]=Y^X$. In such a $\C$ coproducts are
disjoint and $0$ is strict; these are the only extensivity facts we use (see \cite{CLW} for the
theory).

\subsection{Rigidity of the terminal object}

\begin{lemma}\label{lem:E1}
If $\one\cong\one\sqcup Z$ then $Z\cong 0$.
\end{lemma}
\begin{proof}
Let $\theta:\one\sqcup Z\xrightarrow{\ \sim\ }\one$ with injections $i:\one\to\one\sqcup Z$,
$j:Z\to\one\sqcup Z$. Then $\theta i:\one\to\one$ is a map into the terminal object, so
$\theta i=\mathrm{id}_{\one}$. Disjointness makes the pullback of $i$ along $j$ initial; as $\theta$
is invertible, so is the pullback of $\theta i=\mathrm{id}_{\one}$ along $\theta j:Z\to\one$, which is
$Z$ itself. Hence $Z\cong 0$.
\end{proof}

\begin{lemma}\label{lem:E2}
If $\one\not\cong 0$ and $E\cdot\one\cong\one$ for a set $E$, then $|E|=1$.
\end{lemma}
\begin{proof}
$E=\varnothing$ gives $0\cong\one$, excluded. If $|E|\ge 2$, pick $e_0\in E$ and set
$Z:=(E\setminus\{e_0\})\cdot\one$, so $\one\cong\one\sqcup Z$ and Lemma~\ref{lem:E1} gives $Z\cong 0$.
But then $\one$ is a coproduct summand of $Z\cong 0$, giving a map $\one\to 0$; strictness of $0$
forces $\one\cong 0$, excluded.
\end{proof}

\begin{lemma}[Disconnected implies decomposable]\label{lem:E0}
If $\one\not\cong 0$ and $\one$ is not connected, then $\one\cong A\sqcup B$ with $A,B\not\cong 0$.
\end{lemma}
\begin{proof}
Let $(X_d)_{d\in D}$ witness the failure of $\gamma$. Here $\gamma$ is always injective: if
$\iota_{d_0}f=\iota_{d_1}g$ with $f:\one\to X_{d_0}$, $g:\one\to X_{d_1}$ and $d_0\neq d_1$, the
disjointness pullback yields $\one\to 0$, excluded; and if $d_0=d_1$ then $f=g$ because injections are
monic. So $\gamma$ fails to be surjective: some $h:\one\to\coprod_d X_d$ factors through no injection.
Pulling the coproduct decomposition back along $h$, infinitary extensivity gives
$\one\cong\coprod_d A_d$ with $A_d:=h^*X_d$. Not all $A_d$ are initial (else $\one\cong 0$); not
exactly one is non-initial (else $h$ factors through an injection); so at least two are, and grouping
gives $\one\cong A\sqcup B$ with both summands non-initial.
\end{proof}

\subsection{Extensive rigidity}

\begin{theorem}[Extensive rigidity]\label{thm:B}
Let $\C$ be nontrivial, infinitary lextensive and cartesian closed. If $\C$ is $\comp$-admissible,
then $\one$ is connected.
\end{theorem}
\begin{proof}
Suppose $\one$ is not connected. By Lemma~\ref{lem:E0}, $\one\cong A\sqcup B$ with $A,B\not\cong 0$,
and extensivity gives an equivalence $\C\simeq\C/A\times\C/B=:\C_1\times\C_2$; both factors are
infinitary lextensive and nontrivial. Since $\C_1\times\C_2\simeq\C$ is cartesian closed, so is each
factor, and all structure is computed componentwise -- limits, colimits, and exponentials
$Y^X=(Y_1^{X_1},Y_2^{X_2})$; in particular $Y^0\cong\one$ and $Y^{\one}\cong Y$ in each factor. Under
the equivalence $\one\leftrightarrow(\one_{\C_1},\one_{\C_2})$ and
$A\leftrightarrow(\one_{\C_1},0)$. Take any $T$ with $|T|\ge 2$; then
\[
   [A,\,T\cdot\one]\;=\;\big((T\cdot\one_{\C_1})^{\one_{\C_1}},\,(T\cdot\one_{\C_2})^{0}\big)
   \;\cong\;\big(T\cdot\one_{\C_1},\,\one_{\C_2}\big).
\]
By the Shape Lemma there is a set $E$ with $[A,\,T\cdot\one]\cong E\cdot\one=(E\cdot\one_{\C_1},\,
E\cdot\one_{\C_2})$. Componentwise: $\one_{\C_2}\cong E\cdot\one_{\C_2}$ gives $|E|=1$ by
Lemma~\ref{lem:E2} in $\C_2$, and then $T\cdot\one_{\C_1}\cong\one_{\C_1}$ gives $|T|=1$ by
Lemma~\ref{lem:E2} in $\C_1$ -- contradicting $|T|\ge 2$.
\end{proof}

In one sentence: $\Fam(\C^{\op})$ gives an object \emph{one} external shape set, but when $\C$ splits
the composite $\ext p\ext q$ demands one shape set per factor, and $[A,\,T\cdot\one]=(T\cdot\one_1,
\one_2)$ is exactly the demand for the two incompatible counts $|T|$ and $1$ at once.

\begin{corollary}\label{cor:setxset}
$\Setb\times\Setb$ is not $\comp$-admissible: its unit $(1,1)$ is disconnected, and the obstruction
is already visible at $p=\gena{A}$ for $A$ a nontrivial summand of $\one$.
\end{corollary}

\subsection{The forced package}

On the extensive pole, admissibility hands over connectedness, and connectedness hands over the rest.
The following holds over \emph{any} closed base and is the reason the left adjoint is unconditional
in $q$.

\begin{theorem}[Reindexing left adjoint]\label{thm:leftadj}
Let $\C$ be closed symmetric monoidal with small coproducts and connected unit $\one$. Let
$q=(T,Q)$, and suppose $p\comp q$ exists for every $p$. Then
\[
   F_q\;\dashv\;({-})\comp q,\qquad F_q(R,U):=\big(R,\ \rho\mapsto\ext q(U_\rho)\big)=\Fam(\ext q^{\op}),
\]
with no condition on $q$.
\end{theorem}
\begin{proof}
Write $p\comp q=(D,N)$, so $\coprod_d[N_d,X]\cong\ext p\ext q(X)=\coprod_s[P_s,\ext q X]$ naturally.
Fix $r=(R,U)$. Using \eqref{eq:hom}, then $\C(N,U)\cong\C(\one,[N,U])$ (closedness), then $\gamma$
twice (connectedness):
\[
\begin{aligned}
   \Fam(r,\,p\comp q)
   &=\prod_\rho\coprod_{d}\C(N_d,U_\rho)
    =\prod_\rho\coprod_{d}\C(\one,[N_d,U_\rho])\\
   &\cong\prod_\rho\C\Big(\one,\coprod_d[N_d,U_\rho]\Big)
    =\prod_\rho\C\big(\one,\ext p(\ext q U_\rho)\big)\\
   &=\prod_\rho\C\Big(\one,\coprod_s[P_s,\ext q U_\rho]\Big)
    \cong\prod_\rho\coprod_s\C(\one,[P_s,\ext q U_\rho])\\
   &=\prod_\rho\coprod_s\C(P_s,\ext q U_\rho)
    =\Fam\big(F_q r,\,p\big).
\end{aligned}
\]
Every step is a natural isomorphism, so the composite is natural in $r$ and $p$.
\end{proof}

\begin{remark}[One comparison map, two theorems]\label{rem:onelemma}
The proof uses $\gamma$ and nothing else -- no distributivity, and no chosen presentation of
$p\comp q$. It is literally the same map $\gamma$ that Theorem~\ref{thm:T1} inverts for fullness.
Thus \emph{fullness of $\ext{-}$ and existence of the reindexing left adjoint are the same lemma
applied twice}: to the probe families $([Q_t,Z])_t$ for fullness, and to $([N_d,U_\rho])_d$ and
$([P_s,\ext q U_\rho])_s$ for the adjunction. Over $\Setb$, $\gamma=\mathrm{id}$ and the reindexing
adjoint recovers \cite[Prop.~6.57]{NiuSpivak}.
\end{remark}

\begin{corollary}\label{cor:extensive-equiv}
For $\C$ nontrivial infinitary-lextensive cartesian-closed and $\comp$-admissible, the unit $\one$ is
connected, the extension $\ext{-}$ is fully faithful, and $({-})\comp q$ has a left adjoint for every
$q$. This package is one-directional: the converse of Theorem~\ref{thm:B} fails, since connectedness
alone does not imply admissibility (Theorem~\ref{thm:G}). The exact admissibility criterion on this
pole is the subject of Section~\ref{sec:char}.
\end{corollary}

\section{The extensive-pole characterisation}\label{sec:char}

The Shape Lemma is necessary (Lemma~\ref{lem:shape}). On the extensive pole it is also
\emph{sufficient}, and this turns the necessary condition into a complete, checkable characterisation
of admissibility. The engine is an \emph{external distributive law}; the subtlety -- the reason a
connected-unit topos can still fail, as $\Gl((-)^2)$ will in Section~\ref{sec:inadmissible} -- is that
this external law is strictly stronger than the internal distributivity that every topos already
carries.

Throughout this section $\C$ is infinitary lextensive; from Theorem~\ref{thm:char} on, also cartesian
closed. Extensivity supplies, for any map $c:P\to T\cdot\one$ into a copower of the terminal, a
canonical decomposition $P\cong\coprod_{t\in T}P^c_t$, where $P^c_t$ is the pullback of $c$ along the
$t$-th injection $\one\to T\cdot\one$ (coproducts are disjoint and universal, so $\C/(T\cdot\one)\simeq
\prod_t\C/\one$).

\begin{definition}[Shape condition and external distributive law]\label{def:SD}
Write \textbf{(S)} for the standing conclusion of the Shape Lemma: for every $P$ and every set $T$,
$[P,\,T\cdot\one]$ is a copower of $\one$. Write \textbf{(D)} for the \emph{external distributive law}:
for every $P$, every set $T$ and every family $(Y_t)_{t\in T}$, the canonical comparison
\[
   \kappa:\ \coprod_{c\in\C(P,\,T\cdot\one)}\ \prod_{t\in T}[P^c_t,\,Y_t]
   \ \longrightarrow\ \Big[P,\ \coprod_{t\in T}Y_t\Big],
\]
which on the $c$-th summand sends a tuple $(g_t:P^c_t\to Y_t)_t$ to
$P\cong\coprod_t P^c_t\xrightarrow{\coprod g_t}\coprod_t Y_t$, is an isomorphism, naturally in
$(Y_t)$.
\end{definition}

Putting every $Y_t=\one$ collapses (D) to (S), so \emph{(D)}$\Rightarrow$\emph{(S)} is immediate. The
substance is the two harder directions.

\begin{theorem}[The external distributive law is sufficient]\label{thm:Dsuff}
Let $\C$ be infinitary-extensive closed symmetric monoidal and cocomplete, and suppose \emph{(D)}
holds. Then $\C$ is $\comp$-admissible. Explicitly, for $p=(S,P)$ and $q=(T,Q)$ the composite
$\ext p\circ\ext q$ is $\ext{r}$ for
\[
   r\;=\;\coprod_{s\in S} r_s,\qquad
   r_s\;=\;\Big(\,\C(P_s,\,T\cdot\one),\ \big(\textstyle\coprod_{t\in T}P^c_t\otimes Q_t\big)_{c}\,\Big).
\]
\end{theorem}
\begin{proof}
Fix $P$ and $q=(T,Q)$, and put $Y_t:=[Q_t,X]$, functorial in $X$. Then, naturally in $X$,
\[
\begin{aligned}
   [P,\,\ext q X]
   &=\Big[P,\ \coprod_{t}[Q_t,X]\Big]\\
   &\cong\;\coprod_{c\in\C(P,\,T\cdot\one)}\ \prod_{t}\,[P^c_t,\,[Q_t,X]]
     &&\text{(D)}\\
   &\cong\;\coprod_{c}\ \prod_{t}\,[P^c_t\otimes Q_t,\ X]
     &&\text{(tensor--hom)}\\
   &\cong\;\coprod_{c}\ \Big[\,\coprod_{t}(P^c_t\otimes Q_t),\ X\Big]
     &&\big([-,X]\text{ preserves }\coprod\big).
\end{aligned}
\]
Thus $[P,\,\ext q(-)]\cong\ext{r_P}$ with $r_P:=\big(\C(P,T\cdot\one),(\coprod_t P^c_t\otimes Q_t)_c\big)$,
an object of $\Fam(\C^{\op})$ (the shape set is small by local smallness, each position object exists by
cocompleteness). Summing over the shapes of $p$,
$\ext p\ext q(-)=\coprod_s[P_s,\ext q(-)]\cong\coprod_s\ext{r_{P_s}}=\ext{\coprod_s r_{P_s}}$. Every
step but (D) is the closed structure or the universal property of the coproduct.
\end{proof}

Sufficiency needs only a closed monoidal -- not cartesian -- structure, and it constructs $\comp$
explicitly. Its converse over an extensive cartesian-closed base is exactly the Shape condition.

\begin{theorem}[The Shape condition \emph{is} the distributive law]\label{thm:char}
For $\C$ nontrivial infinitary-extensive cartesian closed, \emph{(S)}$\iff$\emph{(D)}.
\end{theorem}
\begin{proof}
$(\Leftarrow)$ is above. For $(\Rightarrow)$ assume (S): $[P,\,T\cdot\one]\cong E\cdot\one$ with
$E:=\C(P,\,T\cdot\one)$, for all $P,T$. Fix $P,T,(Y_t)$ and compare the two sides of (D) through their
representables $\C(X,-)$, naturally in $X$. Three standard facts of an infinitary-extensive
cartesian-closed category are used:
\begin{itemize}[leftmargin=1.5em]
\item[(U)] \emph{Universality:} a map $W\to\coprod_i B_i$ is the same data as a decomposition
$W\cong\coprod_i W_i$ with legs $W_i\to B_i$, so
\[
   \C\Big(W,\coprod_i B_i\Big)\cong\coprod_{\{W\cong\coprod_i W_i\}}\ \prod_i\C(W_i,B_i).
\]
This is what \emph{fails} in an additive category, and is why Theorem~\ref{thm:char} is confined to
the extensive pole.
\item[(P)] \emph{Partitions:} taking $B_i=\one$ in (U) identifies $\C(W,\,T\cdot\one)$ with the set of
$T$-indexed decompositions $W\cong\coprod_{t}W_t$.
\item[(Dist)] $(-)\times P$ is a left adjoint, hence preserves coproducts.
\end{itemize}
\emph{Right-hand side.} With $W_c:=\prod_t[P^c_t,Y_t]$, by (U) then the cartesian-closed adjunction,
\[
   \C\Big(X,\ \coprod_{c\in E}W_c\Big)\cong
   \coprod_{\{X\cong\coprod_{c}X_c\}}\ \prod_{c\in E}\prod_{t}\C(X_c\times P^c_t,\ Y_t).
\]
\emph{Left-hand side.} By the adjunction then (U),
\[
   \C\big(X,\ [P,\textstyle\coprod_t Y_t]\big)\cong\C\big(X\times P,\ \textstyle\coprod_t Y_t\big)\cong
   \coprod_{\{X\times P\cong\coprod_t Z_t\}}\ \prod_{t}\C(Z_t,Y_t).
\]
\emph{The shape indices biject.} By (P), the cartesian-closed adjunction, (S), and (P) again,
\[
\begin{aligned}
   \{X\times P\cong\textstyle\coprod_t Z_t\}
   &=\C(X\times P,\,T\cdot\one)\cong\C(X,[P,\,T\cdot\one])\\
   &=\C(X,\,E\cdot\one)=\{X\cong\textstyle\coprod_{c\in E}X_c\},
\end{aligned}
\]
naturally in $X$. \emph{The fibres match.} Tracing the decomposition $\psi$ named by $\varphi$: the map
$X\times P\to T\cdot\one$ restricts on $X_c\times P$ (for $c\in E=\C(P,T\cdot\one)$) to
$X_c\times P\xrightarrow{\mathrm{pr}}P\xrightarrow{c}T\cdot\one$, and since $c$ decomposes
$P\cong\coprod_t P^c_t$, (Dist) gives $X_c\times P\cong\coprod_t X_c\times P^c_t$. Summing over $c$ with
$X\cong\coprod_c X_c$ yields $Z_t=\coprod_{c\in E}X_c\times P^c_t$, whence
$\prod_t\C(Z_t,Y_t)=\prod_{c\in E}\prod_t\C(X_c\times P^c_t,Y_t)$ -- exactly the right-hand fibre. Both
sides therefore represent the same functor of $X$ (naturally in $(Y_t)$), and Yoneda gives
$[P,\coprod_t Y_t]\cong\coprod_{c\in E}\prod_t[P^c_t,Y_t]$, which is (D).
\end{proof}

The single load-bearing use of (S) is the step $\C(X,[P,T\cdot\one])=\C(X,\,E\cdot\one)$: it converts
the internal shape \emph{object} $[P,T\cdot\one]$ into an external \emph{partition} of $X$ indexed by
the \emph{set} $E$. Without it the shape stays internal -- and the coproduct $\coprod_{c\in E}$ on the
right cannot even be formed over a set. This is precisely the internal/external seam of
Remark~\ref{rem:internal}, and precisely what $\Gl((-)^2)$ violates.

When $\C$ carries a connected-components functor the Shape condition becomes a statement about it. The
following reduction (used again for the gluing base) is the bridge to the multiplicativity criterion.

\begin{proposition}\label{prop:reduction}
Let $\C$ be cocomplete cartesian closed with a components adjunction $\pizero\dashv(-)\cdot\one$ and
$\pizero(\one)\cong 1$. Then \emph{(S)} holds if and only if $\pizero$ preserves finite products, and
$T=2$ already detects failure.
\end{proposition}
\begin{proof}
$[P,\,T\cdot\one]$ represents $X\mapsto\C(X\times P,\,T\cdot\one)=\Setb(\pizero(X\times P),T)$. If
$\pizero$ preserves $(-)\times P$ this is $\Setb(\pizero X\times\pizero P,T)=\C(X,(T^{\pizero P})\cdot
\one)$, so $[P,\,T\cdot\one]\cong(T^{\pizero P})\cdot\one$ is a copower of $\one$. Conversely, if
$[P,\,T\cdot\one]\cong E\cdot\one$, evaluating the representable at $X=\one$ (using $\one\times P\cong
P$, $\pizero\one\cong 1$) gives $E=T^{\pizero P}$, so
$\Setb(\pizero(X\times P),T)\cong\Setb(\pizero X\times\pizero P,T)$ naturally in $X$ for every $T$;
as $2$ cogenerates $\Setb$, the canonical comparison $\pizero(X\times P)\to\pizero X\times\pizero P$
is a bijection.
\end{proof}

\begin{corollary}[Characterisation]\label{cor:char}
For $\C$ nontrivial infinitary-extensive cartesian closed, the following are equivalent:
$\C$ is $\comp$-admissible; the Shape condition \emph{(S)} holds; the external distributive law
\emph{(D)} holds. When moreover $\C$ has a components adjunction $\pizero\dashv(-)\cdot\one$ with
$\pizero(\one)\cong1$ (e.g.\ any Grothendieck topos), these hold if and only if $\pizero$ preserves
finite products.
\end{corollary}
\begin{proof}
Admissible $\Rightarrow$ (S) is the Shape Lemma (Lemma~\ref{lem:shape}); (S)$\Rightarrow$(D) is
Theorem~\ref{thm:char}; (D)$\Rightarrow$ admissible is Theorem~\ref{thm:Dsuff}. The $\pizero$
reformulation is Proposition~\ref{prop:reduction}.
\end{proof}

\begin{corollary}[The two obstructions are one]\label{cor:unif}
On the extensive pole, failure of $\pizero$-multiplicativity is the \emph{single} obstruction to
admissibility. In particular the disconnected-unit failure of $\Setb\times\Setb$ and the
connected-unit failure of $\Gl((-)^2)$ are the same phenomenon, $\pizero(X\times Y)\not\cong\pizero
X\times\pizero Y$. Furthermore $\pizero$-multiplicativity implies a connected unit -- via
Corollary~\ref{cor:char} and Theorem~\ref{thm:B}, admissible $\Rightarrow$ connected -- while the
converse fails ($\Gl$), so connectedness of the unit is strictly weaker than admissibility.
\end{corollary}
\begin{proof}
Only the $\Setb\times\Setb$ identification is new. There $\pizero=(X_1,X_2)\mapsto X_1\sqcup X_2$ (left
adjoint to the diagonal $(-)\cdot\one:T\mapsto(T,T)$), and for the componentwise product
$\pizero\big((X_1,X_2)\times(Y_1,Y_2)\big)=X_1Y_1\sqcup X_2Y_2$, whereas
$\pizero(X)\times\pizero(Y)=(X_1\sqcup X_2)\times(Y_1\sqcup Y_2)$; these differ, so $\pizero$ is not
multiplicative. The rest is Corollary~\ref{cor:char}, Theorem~\ref{thm:B}, and
Corollary~\ref{cor:connected-not-sufficient}.
\end{proof}

\begin{remark}[Internal versus external, again]\label{rem:char-internal}
Every Grothendieck topos supports composition of \emph{internal} (Gambino--Kock) polynomial functors
\cite{GambinoKock}; that is the internal distributive law, and it holds unconditionally. What
Theorem~\ref{thm:char} isolates is the \emph{external} law (D), which governs the external-shape
functors $\coprod_s[P_s,-]$, and which is equivalent to (S). The two coincide exactly when the shape
object $[P,T\cdot\one]$ is a copower of $\one$ -- a discrete object -- and $\pizero$ measures the
discrepancy. This is why admissibility is not automatic on the extensive pole, and why our
connected-unit counterexample $\Gl((-)^2)$ is exactly a base where $[P,T\cdot\one]$ acquires a
non-discrete summand.
\end{remark}

\section{The additive pole: collapse}\label{sec:collapse}

On the opposite extreme every object is copower-tiny, $\comp$ degenerates to $\otimes$, and the
adjunction survives only trivially.

\begin{lemma}\label{lem:collapse-tensor}
If every $P_s$ is copower-tiny then $p\comp q=p\otimes q$; this holds in particular over $\Vecfd$, and
over $\Vecb$ whenever $T$ is finite.
\end{lemma}
\begin{proof}
$[P_s,\ext q X]=[P_s,\coprod_t[Q_t,X]]\cong\coprod_t[P_s,[Q_t,X]]=\coprod_t[P_s\otimes Q_t,X]$ using
copower-tininess and the tensor--hom adjunction; sum over $s$. For finite $T$ over additive $\Vecb$,
$\coprod_t=\prod_t$ and every $[P_s,-]$ preserves finite limits, giving the same identity without
tininess.
\end{proof}

\begin{lemma}[Collapse disconnects the unit]\label{lem:D}
If every object of $\C$ is copower-tiny then $\one$ is not connected.
\end{lemma}
\begin{proof}
Suppose $\one$ were connected. With $T=\{1,2\}$ and $X=\one$, the map
$\mu:\C(\one\sqcup\one,\one)\sqcup\C(\one\sqcup\one,\one)\to\C(\one\sqcup\one,\one\sqcup\one)$,
$(d,f)\mapsto\iota_d\circ f$, is the composite of canonical isomorphisms
\[
   \C(\one\sqcup\one,\one\sqcup\one)\cong\C(\one,[\one\sqcup\one,\,2\cdot\one])
   \cong\C(\one,\,2\cdot[\one\sqcup\one,\one])
   \cong 2\cdot\C(\one,[\one\sqcup\one,\one]),
\]
the last term being $\C(\one\sqcup\one,\one)^{\sqcup 2}$,
where the middle step is copower-tininess of $\one\sqcup\one$ and the next is $\gamma$. Each factor is
post-composition with $\iota_d$, so $\mu$ is a bijection. Hence $\mathrm{id}_{\one\sqcup\one}=
\iota_{d_0}\circ f$ for a unique pair, say $d_0=1$ with $f:\one\sqcup\one\to\one$. Pre-composing with
$\iota_2$ gives $\iota_2=\iota_1\circ(f\circ\iota_2)$. But now $\gamma$ at $2\cdot\one$ sends the two
\emph{distinct} elements $(2,\mathrm{id}_{\one})$ and $(1,f\circ\iota_2)$ of
$\C(\one,\one)\sqcup\C(\one,\one)$ to the same map $\iota_2$, contradicting injectivity.
\end{proof}

No cardinality argument appears above, and this is essential: for $\C=\Vecfd$ over $\mathbb F_2$ with
$\dim=1$ both sides of $\gamma$ have four elements, so a count would pass -- the failure is visible
only in the map. (When $\C$ has a zero object, disconnection is immediate: the zero map lies in every
summand, so $\gamma$ is non-injective.)

\begin{theorem}[Collapse verdict]\label{thm:2}
Let $q=(T,Q)$ with $p\comp q=p\otimes q$ for all $p$ (e.g.\ under Lemma~\ref{lem:collapse-tensor}).
Then $({-})\comp q$ has a left adjoint if and only if $|T|=1$, in which case
$F_q(R,U)=(R,[Q,U_\rho])=\Fam(\ext q^{\op})$.
\end{theorem}
\begin{proof}
\emph{Sufficiency.} For $q=(\{*\},Q)$, tensor--hom gives $\Fam(r,p\otimes q)=\prod_\rho\coprod_s
\C(P_s\otimes Q,U_\rho)=\prod_\rho\coprod_s\C(P_s,[Q,U_\rho])=\Fam((R,[Q,U_\rho]),p)$, naturally.
\emph{Necessity.} A right adjoint preserves the terminal object $\one=\gena{0}$
(Lemma~\ref{lem:terminal}), and $\one\comp q=\one\otimes q=(T,0)$ since $0\otimes Q_t\cong 0$; by
Lemma~\ref{lem:iso}, $(T,0)\cong(\{*\},0)$ iff $|T|=1$. Independently, and using no zero object:
products in $\Fam(\C^{\op})$ are $p\times p'=(S\times S',\,P_s\sqcup P'_{s'})$, so for
$p=p'=\gena{I}$ the shape set of $(p\times p')\otimes q$ is $T$ while that of
$(p\otimes q)\times(p'\otimes q)$ is $T\times T$; preservation forces $|T|=|T|^2$.
\end{proof}

The two necessity arguments differ in robustness. The binary-product argument survives any
re-selection of $\comp$ that keeps shape sets multiplicative and (for finite $T$) already forces
$|T|=1$; the terminal-object argument is convention-sensitive, since for infinite $T$ one has
$|T|=|T|^2$ and $({-})\comp q$ can preserve all binary products while failing to preserve the empty
one. The binding probe is the \emph{empty} limit -- the same phenomenon by which containers preserve
connected but not all limits.

\begin{remark}[The probe method, made explicit]\label{rem:fatalprobe}
The comparison map behind Theorem~\ref{thm:2} is $\gamma$ with the probe $\one$ replaced by an
arbitrary object $B$: $\kappa_{B,Z}:\coprod_t\C(B,[Q_t,Z])\to\C(B,\coprod_t[Q_t,Z])$. The probe
$B=\one$ is connectedness; the \emph{fatal} probe is $B=0$, where the left side has $|T|$ elements and
the right side one, forcing $|T|=1$. One functional, the probe varying, decides both poles.
\end{remark}

\section{Three inadmissible bases}\label{sec:inadmissible}

Off the two poles the Shape Lemma turns into an obstruction. We record three inadmissible bases.
Two of them, $\Setb\times\Setb$ and $\Gl((-)^2)$, lie on the extensive pole and so fall under the
characterisation of Section~\ref{sec:char}: both fail by the \emph{single} mechanism of
$\pizero$-non-multiplicativity (Corollary~\ref{cor:unif}), though they differ on whether the unit is
connected. The third, $\Setb_*$, lies off the pole and fails for an unrelated, polynomial-degree
reason. The sharpest consequence is the gluing case: connectedness of the unit is necessary but not
sufficient.

\subsection{$\Setb\times\Setb$: a disconnected unit}
This is Corollary~\ref{cor:setxset}: the unit is disconnected and the Shape Lemma fails already at
$p=\gena{A}$. It is now an instance of a theorem rather than a computation, and by
Corollary~\ref{cor:unif} the failure is the same $\pizero$-non-multiplicativity that defeats the
gluing base below -- here $\pizero(X_1,X_2)=X_1\sqcup X_2$ does not preserve the componentwise
product.

\subsection{Pointed sets: a polynomial-degree obstruction}
Let $\Setb_*$ be pointed sets under smash, with $I=S^0$, coproduct the wedge $\vee$, and
$0\cong\one=*$ a zero object. Since $\one$ is a zero object, one might expect $\Setb_*$ to behave like
the linear pole; it does not.

\begin{theorem}\label{thm:A}
$\Setb_*$ is not $\comp$-admissible.
\end{theorem}
\begin{proof}
Take $p=\gena{S^0\vee S^0}$ and $q=(\{1,2\},(S^0,S^0))$. Since $[S^0,X]\cong X$ and a pointed map out
of a wedge is a pair of pointed maps, $\ext q X\cong X\vee X$ and $\ext p Y\cong Y\times Y$, so
$\ext p\ext q X\cong(X\vee X)\times(X\vee X)$. Suppose $p\comp q=(D,N)$ existed, so
$\bigvee_{d\in D}[N_d,X]\cong(X\vee X)^2$ naturally. Writing $m=|X|-1$ and $n_d=|N_d|-1$, this is the
polynomial identity, for all $m\ge 0$,
\[
   1+\sum_{d\in D}\big((m+1)^{n_d}-1\big)=(2m+1)^2=4m^2+4m+1. \tag{$*$}
\]
Only finitely many $n_d\ge 1$ can survive (else the left side is infinite at $m=1$), and each
$(m+1)^{n}-1$ has degree $n$ with all coefficients positive, so no cancellation occurs and every
$n_d\in\{1,2\}$. Writing $a=\#\{n_d=1\}$, $b=\#\{n_d=2\}$, $(*)$ reads $am+b(m^2+2m)=4m^2+4m$, whence
$b=4$ and $a+2b=4$, i.e.\ $a=-4$, not a cardinality.
\end{proof}

$\Setb_*$ sits on neither pole: $S^0$ is tiny but $S^0\vee S^0$ is not, and there is no
distributivity or additivity to compensate. Its disconnection is by \emph{non-injectivity} of
$\gamma$ (basepoints are identified), unlike the linear pole's zero-object disconnection. It also
refutes a natural conjecture -- that a zero object suffices for admissibility -- in one line, and
shows the two poles are not jointly exhaustive.

\subsection{Artin gluing: connected but inadmissible}\label{sub:gluing}
The decisive base is the Artin gluing $\Gl((-)^2)=(\Setb\downarrow(-)^2)$, whose objects are directed
graphs $(A,B,\beta)$ -- a vertex set $B$, an edge set $A$, and endpoints $\beta:A\to B\times B$.

\begin{proposition}\label{prop:gl-topos}
$\Gl((-)^2)$ is a Grothendieck topos; in particular cartesian closed, cocomplete, and infinitary
lextensive. Its terminal object $\one$ is a single looped vertex, its product is the tensor
(Kronecker) product of graphs, and $T\cdot\one=(T,T,\mathrm{diag})$ is $T$ looped vertices.
\end{proposition}
\begin{proof}
$(-)^2=\Setb(2,-)$ is right adjoint to $2\times(-)$, hence left exact; Artin gluing of a left-exact
functor between Grothendieck toposes is a Grothendieck topos \cite[A2.1.12]{MacLaneMoerdijk},
\cite{CarboniJohnstone}, and every Grothendieck topos is infinitary extensive. The description of the
finite (co)limits is the standard comma-category computation.
\end{proof}

The base carries two $\Setb$-valued functors that must not be confused. \emph{Points}
$U=\C(\one,-)$ sends a graph to its set of loops $\{a:\beta(a)_0=\beta(a)_1\}$; loops do not cross
components, so $U$ preserves coproducts and \emph{the unit is connected}. \emph{Components}
$\pizero$ sends a graph to its set of connected components; it is left adjoint to $(-)\cdot\one$,
because a morphism $X\to T\cdot\one=(T,T,\mathrm{diag})$ is exactly a map $B\to T$ constant on edges,
i.e.\ $\Hom(X,T\cdot\one)=\Setb(\pizero X,T)$. The reduction of the Shape condition to
$\pizero$-multiplicativity is Proposition~\ref{prop:reduction}.

\begin{theorem}[Gluing obstruction]\label{thm:G}
$\Gl((-)^2)$ is not $\comp$-admissible; the Shape Lemma fails at $P=K$, $T=2$, where $K$ is the single
edge.
\end{theorem}
\begin{proof}
$K$ (vertices $\{0,1\}$, edges $\{0\to1,1\to0\}$) is connected and bipartite, so $\pizero K=1$. By
Weichsel's theorem \cite{Weichsel}, the tensor product of two connected bipartite graphs has exactly
two components, so $\pizero(K\times K)=2\neq1=\pizero K\cdot\pizero K$: $\pizero$ is not
multiplicative, and by Proposition~\ref{prop:reduction} the Shape Lemma fails. Concretely, if
$[K,\,2\cdot\one]$ were a copower of $\one$ it would be $2\cdot\one$ (forced by evaluation at $\one$),
but $\Hom(K,[K,2\cdot\one])=\Hom(K\times K,2\cdot\one)=\Setb(\pizero(K\times K),2)=4\neq2$. Indeed
$[K,\,2\cdot\one]\cong 2\cdot\one\sqcup K$: two looped vertices carrying the two points, plus a
\emph{bald edge} -- a component with no loop, hence no point -- which is exactly the summand
obstructing the Shape Lemma.
\end{proof}

\begin{corollary}[Connected is not sufficient]\label{cor:connected-not-sufficient}
Connectedness of the unit is necessary for $\comp$-admissibility on the extensive pole
(Theorem~\ref{thm:B}) but not sufficient: $\Gl((-)^2)$ has a connected unit yet is inadmissible.
\end{corollary}

Thus the extensive pole itself splits, exactly along the characterisation of
Corollary~\ref{cor:char}. Where $\pizero$ is multiplicative (e.g.\ $\Setb$, where
$\pizero=\mathrm{Id}$) the base is admissible; where pairing two connected objects can disconnect
them, it is not. This does \emph{not} weaken the sufficiency of the Shape Lemma -- $\Gl$ is
inadmissible precisely because it \emph{fails} the Shape Lemma, $[K,2\cdot\one]$ not being a copower
of $\one$ -- but it shows the weaker hypothesis of a merely connected unit is not enough: what
admissibility needs is the full Shape condition, equivalently $\pizero$-multiplicativity.

\begin{remark}[Internal versus external]\label{rem:internal}
The topos $\Gl((-)^2)$ \emph{does} carry composition of internal (Gambino--Kock) polynomial functors
\cite{GambinoKock}; what fails is closure of the external-shape functors $\coprod_s X^{P_s}$. The
composite $\ext{\gena{K}}\circ\ext{(2,(0,0))}$ equals $[K,2\cdot\one]\cong 2\cdot\one\sqcup K$, which
needs a shape of size $3$ with one position object $K\not\cong\one$: a genuine object of
$\Fam(\Gl^{\op})$, but not the copower
of $\one$ that admissibility would force. Theorem~\ref{thm:G} is precisely this external/internal gap,
localised to a base where it bites.
\end{remark}

\section{The middle region is inhabited}\label{sec:middle}

Extensive rigidity (Theorem~\ref{thm:B}) is a \emph{sufficient} condition for connectedness, not a
characterisation of admissibility. We show its hypotheses are essential by exhibiting an admissible
base with a disconnected unit that is neither pole.

Let $\C=\Setb\times\Vecfd$ over a field $k$, with componentwise structure: $\otimes=(\times,\otimes_k)$,
unit $I=(1,k)$, internal hom $[X,Y]=(Y_s^{X_s},\Hom_k(X_v,Y_v))$, and coproduct
$X\sqcup Y=(X_s\sqcup Y_s,\ X_v\oplus Y_v)$ (the $\Vecfd$-factor is the biproduct). Objects of
$\Fam(\C^{\op})$ are $p=(S,(A_s,V_s))$ with $A_s$ a set and $V_s$ finite-dimensional, and
\[
   \ext p(X)=\Big(\coprod_{s} X_s^{A_s},\ \bigoplus_s\Hom_k(V_s,X_v)\Big).
\]

\begin{proposition}[Coordinates]\label{prop:coords}
$\C=\Setb\times\Vecfd$ is non-collapse, non-cartesian, and has a disconnected unit.
\end{proposition}
\begin{proof}
\emph{Non-collapse}: $(A,V)$ is copower-tiny iff $(-)^A$ preserves $\sqcup$ on the $\Setb$-factor,
i.e.\ $|A|\le 1$; objects with $|A|\ge2$ are not tiny. \emph{Non-cartesian}: the categorical product
has $\Vecfd$-part $X_v\oplus Y_v$ (the biproduct), and $\oplus_k\neq\otimes_k$, so $\times\neq\otimes$.
\emph{Disconnected unit}: $\C(I,-)=(-)_s\times|(-)_v|$ where $|-|:\Vecfd\to\Setb$ is the underlying-set
functor, which does not preserve $\oplus$ (a clean witness: $\C(I,(2,k^2))$ has $2|k|^2$ elements
while $\coprod_{i=1,2}\C(I,I)$ has $2|k|$, unequal for $|k|\ge2$).
\end{proof}

\begin{theorem}\label{thm:setxvec}
On the locus $\mathcal P\subseteq\Fam(\C^{\op})$ of objects with finite \emph{vec-support}
$\{s:V_s\neq0\}$, the base $\Setb\times\Vecfd$ is $\comp$-admissible: for all $p,q\in\mathcal P$ there
is $r=p\comp q\in\mathcal P$ and an explicit natural isomorphism $\ext r\cong\ext p\circ\ext q$.
Consequently the hypotheses of Theorem~\ref{thm:B} are essential -- admissibility does not imply
connectedness.
\end{theorem}
\begin{proof}[Proof sketch]
The locus is forced: the $\Vecfd$-part of $\ext p(X)$ is $\bigoplus_s[V_s,X_v]$, finite-dimensional for
all fd $X_v$ iff the vec-support is finite (Lemma~\ref{lem:1.1} below). On $\mathcal P$ the two factors
compose independently. The $\Setb$-part is the classical container composite, with shape
$D=\coprod_s T^{A_s}$ and positions $C_{(s,f)}=\coprod_{a\in A_s}B_{f(a)}$; the bijection
$(s,f,g)\mapsto(s,h)$, $h(a)=(f(a),g|_{B_{f(a)}})$, is a natural isomorphism. The $\Vecfd$-part, using
that each $V_s$ is copower-tiny in $\Vecb$ and the sums are finite, is
\[
   \bigoplus_s[V_s,\textstyle\bigoplus_t[W_t,X_v]]=\bigoplus_{s,t}[V_s\otimes W_t,X_v],
\]
a finite biproduct, which we realise by \emph{concentrating} it at a single chosen slot $d_0\in D$:
$U_{d_0}=\bigoplus_{s,t}V_s\otimes W_t$, $U_d=0$ otherwise. The comparison is the composite of
tensor--hom and regrouping isomorphisms, natural in $X_v$. Both components are natural isomorphisms and
a morphism of $\Setb\times\Vecfd$ is a pair, so the product is a natural isomorphism $\ext r\cong
\ext p\ext q$; and $r\in\mathcal P$ since its vec-support is $\subseteq\{d_0\}$. As $\C$ is
non-cartesian with disconnected unit (Proposition~\ref{prop:coords}), Theorem~\ref{thm:B} does not
apply, and no contradiction arises.
\end{proof}

\begin{lemma}\label{lem:1.1}
$\ext p$ restricts to an endofunctor of $\Setb\times\Vecfd$ iff $p$ has finite vec-support.
\end{lemma}
\begin{proof}
$\bigoplus_s[V_s,X_v]$ is finite-dimensional for all fd $X_v$ iff only finitely many $[V_s,X_v]$ are
nonzero, i.e.\ iff $\{s:V_s\neq0\}$ is finite; the $\Setb$-part is a set for any $S,A_s$.
\end{proof}

The mechanism is a \emph{factoring} of two sources of admissibility.

\begin{proposition}[Admissibility is per-object]\label{prop:absorptive}
Call $P\in\C$ \emph{absorptive} if $X\mapsto[P,\ext q X]$ is an extension for every $q$. Then $\C$ is
$\comp$-admissible iff every object of $\C$ is absorptive. There are two sources of absorptivity: a
copower-tiny (\emph{flexible}) $P$ gives $[P,\ext q X]=\coprod_t[P\otimes Q_t,X]$; and a
distributive (\emph{rigid}) $P$ gives $[P,\coprod_t Y_t]=\coprod_{f:P\to T}\prod_a Y_{f(a)}$.
\end{proposition}
\begin{proof}
$\ext p\ext q X=\coprod_s[P_s,\ext q X]$: a coproduct of extensions is an extension, so if every $P_s$
is absorptive then $\C$ is admissible; conversely take $p=\gena{P}$. The two formulas are the tiny and
distributive computations.
\end{proof}

In $\Setb\times\Vecfd$ an object $(A,V)$ is absorptive because $A$ is rigid (the $\Setb$-factor is
distributive) and $V$ is flexible ($\Vecfd$ is all-tiny). This is exactly why $\Setb\times\Setb$, with
\emph{both} factors rigid but a disconnected unit, is inadmissible while $\Setb\times\Vecfd$ is not:
the operative dichotomy is rigid versus flexible, not cartesian versus non-cartesian. We record the
resulting picture as a conjecture.

\begin{conjecture}[Decomposition]\label{conj:decomp}
Over a closed symmetric monoidal cocomplete base, every absorptive object is a tensor of a
copower-tiny part and a distributive part; consequently every ``nice'' admissible base decomposes as
\emph{rigid-extensive $\times$ flexible-additive}, and every ``nice'' inhabitant of the middle region
so decomposes. If true, no \emph{irreducible} (indecomposable) inhabitant exists among nice bases.
\end{conjecture}

\begin{remark}[The open irreducible case]\label{rem:gap1}
Whether an \emph{irreducible} middle inhabitant exists is open. Such a base would be indecomposable,
admissible, non-collapse, with disconnected unit. By Proposition~\ref{prop:absorptive} it is not any
$\mathrm{fg}$-module category (all such are all-tiny, hence collapse). If additive it needs a non-tiny
object and infinite direct sums -- exactly the open case of $\comp$ on $\Fam(\Vecb^{\op})$ with
infinite shape sets. If non-additive its unit is disconnected by non-injectivity of $\gamma$, and the
only such base known, $\Setb_*$, is inadmissible (Theorem~\ref{thm:A}). The irreducible case is thus
pinned between two identified open doors.
\end{remark}

\section{Why the question has this shape}\label{sec:determinacy}

The map assembled above -- rigid extensive pole, flexible additive pole, thin inhabited middle,
three inadmissible bases -- is organised by a single fact about the extension. On the extensive pole,
admissibility already forces the extension to be injective on objects; off it, it need not be, and
then $\comp$ is not even determined by $\C$.

\begin{theorem}[Determinacy]\label{thm:determinacy}
\begin{enumerate}[leftmargin=2em]
\item If $\one$ is connected then $\ext{-}$ is fully faithful (Theorem~\ref{thm:T1}), hence injective
on objects up to isomorphism; so for admissible $\C$ the object $p\comp q$ is determined up to
canonical isomorphism by $\ext{p\comp q}\cong\ext p\ext q$, and ``$({-})\comp q$ has a left adjoint''
is a property of $\C$.
\item If $\one$ is disconnected this can fail. Over $\Vecfd$, the objects $(\{*\},k^2)$ and
$(\{1,2\},k)$ both present $X\mapsto X\oplus X$ but are not isomorphic (Lemma~\ref{lem:iso}, shape sets
of sizes $1$ and $2$); so $p\comp q$ is not pinned down, and $\comp$ is a \emph{choice}.
\end{enumerate}
\end{theorem}
\begin{proof}
(1) A fully faithful functor reflects isomorphisms and is injective on objects up to isomorphism.
(2) $\ext{(\{*\},k^2)}X=[k^2,X]\cong X\oplus X$ and $\ext{(\{1,2\},k)}X=[k,X]\oplus[k,X]\cong X\oplus
X$.
\end{proof}

Determinacy explains every feature of the map. On the extensive pole (Section~\ref{sec:extensive}),
$\comp$ is canonical, so admissibility hands over connectedness, full faithfulness, and the reindexing
adjoint (Corollary~\ref{cor:extensive-equiv}), while admissibility itself is pinned exactly by
$\pizero$-multiplicativity (Corollary~\ref{cor:char}). In the middle region (Section~\ref{sec:middle}), the
disconnected unit is exactly why the construction of $\comp$ on $\Setb\times\Vecfd$ had to
\emph{choose} a concentration slot $d_0$: there is no canonical element of $\coprod_s T^{A_s}$, so
$(p\comp q)\comp r$ and $p\comp(q\comp r)$ are different presentations of the same functor and $\comp$
has no canonical monoidal structure -- only a choice-dependent one, living canonically only in the
image $\ext{\mathcal P}\subseteq[\C,\C]$. On the collapse pole, the same non-faithfulness is why the
convention $\comp:=\otimes$ cannot be discharged, and why (Theorem~\ref{thm:2}) the binary-product
necessity argument -- immune to re-choice -- is load-bearing while the terminal-object one is not.

In short: \emph{the hypothesis that makes the reindexing adjoint exist is the same hypothesis that
makes the base-change question well-posed}. Exactly when $\one$ is disconnected, there is nothing left
for ``the'' composition product to be until one names a $\comp$; and for the two conventions in actual
use -- container substitution over $\Setb$, and $\otimes$ over the additive pole -- the answers to the
adjoint question are respectively ``always'' and ``only for a single shape'', matching connectedness
exactly.

\section{Conclusion}\label{sec:conclusion}

Whether $\Fam(\C^{\op})$ carries the composition product is not a graded question. A single condition,
the Shape Lemma, is necessary in general and -- this is the paper's main quantitative result -- also
sufficient on the extensive pole, where it becomes a complete characterisation: $\Fam(\C^{\op})$
carries $\comp$ if and only if the connected-components functor $\pizero$ preserves finite products.
This forces a near-dichotomy. On the extensive pole admissibility forces a
connected unit, hence a fully faithful extension, a canonical $\comp$, and the reindexing left adjoint
for every $q$; on the additive pole $\comp$ collapses to $\otimes$ and the adjoint survives only for a
single shape. The two poles do not exhaust the admissible bases -- $\Setb\times\Vecfd$ inhabits the
middle by combining a rigid and a flexible factor -- but they bound it, and the characterisation shows
the two extensive obstructions are one: connectedness of the unit, which admissibility yields, does
not conversely suffice, because the Artin gluing $\Gl((-)^2)$ is a connected-unit topos whose
$\pizero$ is still not multiplicative. Behind all of this is Determinacy: a connected unit makes
$\comp$ a property of $\C$, a disconnected one makes it a choice.

\subsection*{A design principle}
For the applications that motivate the polynomial calculus -- process substitution, plug-in
composition, resource semantics -- the map reads as a constraint on the resource base $\C$. A
compositional substitution calculus is available only when $\C$ is set-like enough for a connected
unit \emph{and} its notion of joint resource is ``cartesian-like'' rather than ``tensor-like'': its
connected-components functor must be multiplicative under pairing. A base where combining two
connected resources can disconnect them admits no external substitution operation, even when it is a
topos with a connected unit. By Corollary~\ref{cor:char} this test is not merely necessary but
exactly complete on the cartesian side: test $\pizero(X\times Y)\cong\pizero X\times\pizero Y$, and
the external composition operator exists precisely when the test passes.

\subsection*{Open problems}
The extensive-pole question -- \emph{is the Shape Lemma sufficient?} -- is settled affirmatively by
Corollary~\ref{cor:char}: it is, and the exact content is the \emph{external} distributive law (D),
not the internal distributivity every topos already has. What remains open lies off that pole.
\begin{enumerate}[leftmargin=2em]
\item \emph{A global criterion.} Corollary~\ref{cor:char} characterises admissibility by
$\pizero$-multiplicativity on the extensive pole; the additive pole is admissible instead by
copower-tininess (Lemma~\ref{lem:collapse-tensor}), a different mechanism with no $\pizero$ available
($\one\cong0$). Is there a single criterion covering both poles and the mixed region between them?
Conjecture~\ref{conj:decomp} is one precise form.
\item \emph{Is there an irreducible middle inhabitant?} Remark~\ref{rem:gap1} pins this between the
open additive case ($\Fam(\Vecb^{\op})$, infinite shape sets and infinite-dimensional positions) and
the conjecture that every $\gamma$-non-injective closed base is inadmissible -- a strengthening of
Theorem~\ref{thm:A}.
\item \emph{Does $\comp$ extend to full $\Vecb$ for infinite shape sets?} A non-tiny position breaks
the collapse presentation; whether $X\mapsto[P,\bigoplus_t[Q_t,X]]$ is ever an extension there is
unsettled.
\item \emph{Decomposition.} Prove or refute Conjecture~\ref{conj:decomp}.
\end{enumerate}

\subsection*{Relation to generalised polynomials}
Dorta, Jarvis and Niu \cite{DJN} construct a Dirichlet tensor and a composition product over a general
base and leave base conditions for future work; the characterisation of Section~\ref{sec:char}
(Corollary~\ref{cor:char}) answers, for the external-shape category, a question in that direction. Their composition product is a confirming
instance rather than a competitor: it is defined on a free product completion with no coproducts, so
the hypotheses of Theorem~\ref{thm:B} never apply while its conclusion always holds. Their product is
also \emph{weighted} and differs from ours off the point base (Remark~\ref{rem:djn}); their Dirichlet
tensor agrees with ours.

\begin{remark}[$\comp$ versus $\compDJN$]\label{rem:djn}
Under the identification of the two categories, our composition-representing product $\comp$ and the
weighted product $\compDJN$ of \cite{DJN} have the same shapes and direction-indexing but different
direction objects: at a position they multiply in an extra outer weight, so they agree exactly when
$\C=1$ and diverge already at $\C=2$. Our defining property $\ext{p\comp q}\cong\ext p\circ\ext q$ is
moreover unavailable on their category for $\C\neq1$, since their embedding does not land in
endofunctors. We therefore keep the symbols distinct throughout.
\end{remark}

\begin{thebibliography}{99}

\bibitem{AAG}
M.~Abbott, T.~Altenkirch, and N.~Ghani,
\emph{Containers: constructing strictly positive types},
Theoretical Computer Science \textbf{342} (2005), 3--27
(preliminary version, \emph{Categories of containers}, FoSSaCS 2003).

\bibitem{CLW}
A.~Carboni, S.~Lack, and R.~F.~C.~Walters,
\emph{Introduction to extensive and distributive categories},
Journal of Pure and Applied Algebra \textbf{84} (1993), 145--158.

\bibitem{CarboniJohnstone}
A.~Carboni and P.~T.~Johnstone,
\emph{Connected limits, familial representability and Artin glueing},
Theory and Applications of Categories (1995).

\bibitem{DJN}
J.~Dorta, P.~Jarvis, and N.~Niu,
\emph{Monoidal structures on generalized polynomial categories},
arXiv:2305.05655 (2023).

\bibitem{GambinoKock}
N.~Gambino and J.~Kock,
\emph{Polynomial functors and polynomial monads},
Mathematical Proceedings of the Cambridge Philosophical Society \textbf{154} (2013), 153--192;
arXiv:0906.4931.

\bibitem{MacLaneMoerdijk}
S.~Mac Lane and I.~Moerdijk,
\emph{Sheaves in Geometry and Logic: A First Introduction to Topos Theory},
Springer, 1992.

\bibitem{NiuSpivak}
N.~Niu and D.~I.~Spivak,
\emph{Polynomial Functors: A Mathematical Theory of Interaction},
arXiv:2312.00990 (2023).

\bibitem{Weichsel}
P.~M.~Weichsel,
\emph{The Kronecker product of graphs},
Proceedings of the American Mathematical Society \textbf{13} (1962), 47--52.

\bibitem{proofT1}
MacBeth,
\emph{The fullness criterion for the container extension over a closed base},
Kodamai research note, 2026.

\bibitem{proofconv}
MacBeth,
\emph{Left and right adjoints of the composition product over a base},
Kodamai research note, 2026.

\end{thebibliography}

\end{document}
